% LSPV-5 REPORT::PRIMELQP1::2026-10-02T08:39:00-07:00
\documentclass[11pt]{article}
\usepackage[margin=.9in]{geometry}\usepackage{amsmath,amssymb}
\begin{document}
\section*{PRIME-LQP-1: prime-factor provenance as a SAME\_REALIZATION candidate}
\textbf{STATUS: CANDIDATE with exact negative baseline. P versus NP remains OPEN.}

\paragraph{Motivation.}
Test the proposed use of primes as a multisign/provenance coordinate for successive Boolean operations. This branch is separate from BITRANK1/product-span and does not replace it.

\paragraph{Literature correspondence.}
Prime-power sequence encodings are a classical exact syntax/proof encoding device (Goedel numbering). Provenance semirings independently formalize annotations that retain derivational information beyond extensional output. These support the legitimacy of a provenance layer, but neither supplies a SAT lower bound or a nonrelativizing theorem.

\paragraph{Definition: positional prime provenance.}
For a finite computation history $h=(e_1,\ldots,e_T)$, let $p_i$ be the $i$th prime and let $c(e_i)\ge1$ encode the declared event type (operator, branch, source label, etc.). Define
\[
P(h)=\prod_{i=1}^T p_i^{c(e_i)}.
\]
Unique factorization makes $P$ injective on the encoded event sequence, provided the coding map and length convention are fixed. Operator-type-only products without positional primes are insufficient because commutativity loses order.

\paragraph{PROVED/EXACT negative baseline: semantic contexts cannot distinguish equivalent histories.}
If Boolean expressions $F,G$ compute the same Boolean function, then for every ordinary Boolean context $C[\cdot]$,
\[
F\equiv G\Longrightarrow C[F]\equiv C[G].
\]
Therefore $P(F)\ne P(G)$ cannot make an ordinary Boolean continuation distinguish $F$ from $G$. Any consumer that reads $P$ is using an enriched provenance semantics, not ordinary extensional Boolean semantics.

\paragraph{PROVED/EXACT negative baseline: deterministic trace-prime labels are reacquirable.}
Fix deterministic polynomial-time $A$ and a polynomial-time transducer that, from $(x,\operatorname{Tr}_A(x))$, emits the prime-event sequence or any polynomial-length residue/digest of it. Then the annotation is in FP by rerunning $A$ and the transducer. This is an instance of the existing TRACECLOSE1/REACQ1 closure. Hence prime labels attached only to an ordinary deterministic polynomial trace cannot be the final nonrelativizing hinge.

\paragraph{Representation warning.}
The integer $P(h)$ itself may have polynomial bit length for polynomial $T$ and polynomial-size exponents, but explicitly factoring an arbitrary supplied integer is not the relevant cost: the generator already knows the factors. A lower bound may not be obtained by hiding a known trace behind a difficult representation conversion. Residues $P\bmod q_j$ are likewise polynomial-time forward-computable when the trace is known.

\paragraph{The only surviving candidate.}
Bind event primes to a SAME\_REALIZATION object that is not determined by the current semantic state alone:
\[
p_{(\mathrm{source},\mathrm{witness},\mathrm{gate},\mathrm{branch},\mathrm{position})}.
\]
A downstream consumer must require a property of the \emph{same} bound realization, not merely inspect an artificial tag. The required theorem would have form
\[
\text{valid downstream object}\Rightarrow
\text{recover/verify nontrivial information about the same prime-bound history},
\]
together with a proof that this information cannot be regenerated, replaced, canonicalized, or supplied by the maximal semantic/algebraic oracle within the charged model.

\paragraph{PRIME-LQP experiment design.}
For Boolean expressions/circuits through $n=5$:
\begin{enumerate}
\item enumerate paired histories with identical truth tables but distinct positional-prime provenance;
\item ordinary-continuation control: verify no Boolean context distinguishes equivalent histories;
\item provenance-aware control: allow direct reading of $P$ and mark resulting distinction TRIVIAL/ARTIFICIAL;
\item SAME\_REALIZATION test: expose only a declared residue vector $(P\bmod q_1,\ldots,P\bmod q_s)$ or selected divisibility predicates tied to legal source/witness events;
\item maximal-oracle attack: supply every semantic answer and test whether the downstream legal consumer still requires a prime-bound fact unavailable from semantics;
\item RECOMP attack: rerun/re-encode the deterministic source and determine whether the prime-bound fact is reacquired in polynomial time;
\item n=5 mandatory regression: reject any invariant that passed n=4 but collapses at n=5.
\end{enumerate}

\paragraph{Candidate observable.}
For two histories $h_1,h_2$, $\gcd(P(h_1),P(h_2))$ exactly identifies common positional events only when the prime assignment itself uses the same position/event namespace. More robustly, use a set of event-specific primes and track divisibility. This is an exact bookkeeping identity, not hardness.

\paragraph{Barrier audit.}
The construction is syntactic and normally relativizes. A semantic oracle can ignore annotations unless the model independently requires them. An algebraic oracle can also be granted the annotation map. Goedel-style injectivity is not nonalgebrization. Prime factorization hardness is inadmissible as an unconditional assumption and is unnecessary for forward annotation. DEAD/COPY/DELAY change provenance without changing semantics; therefore any proposed cost based on raw prime mass is invalid. INTERP/RECOMP directly attacks acquisition. Horn, 2-SAT, affine/Tseitin, bounded-treewidth, padding, renaming, and independent-product controls must remain easy.

\paragraph{Relation to LQP.}
Let $Q_{\rm prov}$ be the prime-annotated realization space and $Q_{\rm sem}$ its Boolean semantic quotient. The map
\[
\pi:Q_{\rm prov}\to Q_{\rm sem}
\]
is many-to-one. PRIME-LQP asks whether a source-legal downstream relation is well-defined only on $Q_{\rm prov}$, rather than factoring through $\pi$. If it factors through $\pi$, provenance is irrelevant. If it does not, one still needs an arbitrary-algorithm theorem forcing the solver to realize that enriched relation.

\paragraph{Corpus reconciliation.}
BITRANK1 was read back at blob 0d70bd133287795326e8a90e65b0434d4c734373 with OPEN status and verification sentinel; its characteristic-two rank audit therefore predates and remains independent of this mod-q/prime discussion. Concurrent LSPV-2 EXCESS1 and LSPV-3 RAMCOMP1 commits were observed and are not overwritten here.

\paragraph{Single next target.}
Construct the smallest source-matched PRIME-LQP pair $(h_1,h_2)$ with identical current SU/vMS/RMV semantic state and different source/witness-bound prime provenance, then search automatically for a later \emph{legal} SAME\_REALIZATION restriction that distinguishes them while denying the consumer direct access to an artificial provenance tag. If no such legal restriction exists through the n=5 battery, downgrade prime provenance to bookkeeping only.

\begin{thebibliography}{9}
\bibitem{GKT07} T. J. Green, G. Karvounarakis, V. Tannen, Provenance Semirings, PODS 2007, 31--40.
\bibitem{GodelCoding} Standard prime-power Goedel coding of finite symbol sequences; see Stanford Encyclopedia of Philosophy, supplement on Goedel numbering.
\end{thebibliography}
\end{document}
% END_OF_VERIFIED_REPORT::LSPV-5::PRIMELQP1
